\section{Kiến trúc hệ thống}

\subsection{Bản đồ codebase}

\begin{lstlisting}[style=rae]
configs/
  stage1/
  stage2/
src/
  stage1/
  stage2/
  utils/
  train.py
  sample.py
  sample_ddp.py
  stage1_sample.py
  stage1_sample_ddp.py
  build_fid_stats.py
  evaluate_fid.py
src_jax/
  vendor.py
  config_adapter.py
  stage2_runtime.py
  stage1_runtime.py
  hf_utils.py
  train.py
  sample.py
  sample_ddp.py
  stage1_sample.py
  push_hf.py
\end{lstlisting}

\subsection{Stage 1: Representation Autoencoder}

Thành phần gốc của Stage 1 nằm trong \texttt{src/stage1/rae.py}. Lớp
\texttt{RAE} kết hợp:

\begin{itemize}[leftmargin=1.5em]
  \item encoder ở \texttt{src/stage1/encoders/},
  \item decoder ở \texttt{src/stage1/decoders/},
  \item latent normalization qua \texttt{normalization\_stat\_path},
  \item latent noising qua \texttt{noise\_tau}.
\end{itemize}

Hành vi chính:

\begin{itemize}[leftmargin=1.5em]
  \item \texttt{encode(x)} resize đầu vào về đúng độ phân giải của encoder,
  chuẩn hóa theo image processor, chạy encoder, reshape latent về dạng 2D nếu
  cần, rồi áp normalization;
  \item \texttt{decode(z)} đảo normalization, đổi latent về dạng sequence nếu
  cần, rồi giải mã thành ảnh RGB bằng decoder ViT.
\end{itemize}

Ở đường JAX, \texttt{src\_jax/stage1\_runtime.py} không tự viết lại toàn bộ RAE
trong repo này. Nó gọi backend NNX, nhưng vẫn dùng đúng các giá trị cấu hình
trong block \texttt{stage\_1}.

\subsection{Stage 2: DiT trong không gian latent}

Implementation PyTorch/XLA chính nằm trong \texttt{src/stage2/models/DDT.py}.
Mô hình xuất hiện nhiều nhất trong repo này là \texttt{DiTwDDTHead}.

Các đặc điểm quan trọng:

\begin{itemize}[leftmargin=1.5em]
  \item đầu vào là latent thay vì pixel;
  \item embedding thời gian bằng Gaussian Fourier features;
  \item class conditioning bằng \texttt{LabelEmbedder};
  \item tùy chọn ROPE, RMSNorm, SwiGLU, qk-norm;
  \item hỗ trợ classifier-free guidance và autoguidance khi sampling.
\end{itemize}

Đường JAX ánh xạ \texttt{stage\_2.target} sang backbone tương ứng của backend
NNX:

\begin{itemize}[leftmargin=1.5em]
  \item \texttt{DiTwDDTHead} $\rightarrow$ \texttt{lightning\_ddt},
  \item \texttt{LightningDiT} $\rightarrow$ \texttt{lightning\_dit},
  \item các target kiểu DiT khác $\rightarrow$ \texttt{dit}.
\end{itemize}

\subsection{Transport objective}

Phần này nằm trong \texttt{src/stage2/transport/transport.py}. Đây là lớp trung
gian giữa latent diffusion model và train loop.

Nhiệm vụ của transport:

\begin{itemize}[leftmargin=1.5em]
  \item lấy mẫu thời gian \texttt{t};
  \item xây dựng đường nối giữa noise và dữ liệu thật;
  \item tính target phù hợp với loại output của model;
  \item trả về drift function để phục vụ ODE hoặc SDE sampler.
\end{itemize}

Ở đường JAX, adapter không giữ lại implementation transport gốc này. Nó ánh xạ
các ý nghĩa tương đương sang interface của backend NNX, chủ yếu là SiT với
sampling kiểu Euler.

\subsection{Lớp adapter JAX}

\begin{center}
\begin{longtable}{>{\raggedright\arraybackslash}p{0.28\linewidth} >{\raggedright\arraybackslash}p{0.62\linewidth}}
\toprule
File & Vai trò \\
\midrule
\endhead
\texttt{src\_jax/vendor.py} & Bootstrap \texttt{diffuse\_nnx} vào \texttt{\~{}/.cache/rae\_jax/diffuse\_nnx} và pin commit \\
\texttt{src\_jax/config\_adapter.py} & Chuyển YAML OmegaConf của repo sang config backend NNX \\
\texttt{src\_jax/stage2\_runtime.py} & Train, sampling, load checkpoint PyTorch hoặc Orbax, FID glue, wandb bridge \\
\texttt{src\_jax/stage1\_runtime.py} & Reconstruction Stage 1 trên đường JAX \\
\texttt{src\_jax/hf\_utils.py} & Upload file hoặc thư mục lên Hugging Face \\
\bottomrule
\end{longtable}
\end{center}

\subsection{Hai mô hình thực thi}

\subsubsection{TorchXLA}

Các entrypoint distributed trong \texttt{src/} dùng \texttt{xmp.spawn(...)} và
mỗi TPU core chạy một process riêng. Các primitive quan trọng:

\begin{itemize}[leftmargin=1.5em]
  \item \texttt{DistributedSampler},
  \item \texttt{ParallelLoader},
  \item \texttt{xm.optimizer\_step(...)},
  \item \texttt{xm.rendezvous(...)} và \texttt{xm.mesh\_reduce(...)}.
\end{itemize}

\subsubsection{JAX / NNX}

Đường \texttt{src\_jax/} không tự quản lý toàn bộ graph, checkpoint và sampler ở
mức thấp. Nó dựa vào backend NNX để:

\begin{itemize}[leftmargin=1.5em]
  \item tạo model, optimizer, sampler và EMA,
  \item quản lý checkpoint Orbax,
  \item chạy sampling và FID trong định dạng tensor NHWC của JAX,
  \item mở rộng sang nhiều process JAX nếu môi trường đã cấu hình phù hợp.
\end{itemize}

\subsection{Checkpoint và tương thích trọng số}

Một điểm quan trọng của nhánh JAX là vẫn tận dụng được checkpoint cũ:

\begin{itemize}[leftmargin=1.5em]
  \item nếu \texttt{stage\_2.ckpt} là file PyTorch \texttt{.pt}, adapter sẽ
  port state dict sang NNX để inference hoặc khởi tạo train;
  \item nếu \texttt{stage\_2.ckpt} là thư mục checkpoint Orbax, adapter sẽ
  restore trực tiếp run JAX trước đó;
  \item decoder Stage 1 ở đường JAX vẫn dùng được checkpoint decoder PyTorch.
\end{itemize}
